\section{Statistical protocol}\label{sec:met-stats}
Every powered closed-loop cell is replicated over $n{=}30$ seeds per topology; the exploratory sweep that preceded them ran three seeds per cell and is labelled as such throughout.
The test for a delay effect is a one-sample $t$-test on the per-seed
relative change against each baseline, computed on the same quantity as the
reported CI, and I pool across topologies ($n{=}90$) where the hypothesis
is topology-independent.
The explanatory study's ANOVA uses the project's
pure-mathematics regularised incomplete beta, so no library version can
silently change it. The closed-loop $t$ and Wilcoxon $p$-values are SciPy's,
recomputable from the committed per-seed JSON. For the factorial explanatory analysis I report effect sizes
(partial $\eta^2$) alongside $F$ statistics, and for headline comparisons I state
direction, magnitude and $p$ together rather than significance alone.

Two commitments were made in advance. First, the honest-null policy: a
non-significant result at $n{=}30$ is reported with its magnitude and
uncertainty, not suppressed or re-run until significant. Second, no
mid-experiment stopping: cell counts were fixed by the factorial design
before launch, and the resumable runner (\S\ref{sec:imp-harness}) re-runs
only cells that produced no measurement (crash, probe failure), never cells
whose answer was unwelcome.

